% !TEX root = ../main.tex

\section{Getting Started with NeoSetzer}
\label{sec:getting-started}

This folder is your own copy of the NeoSetzer example project. It is safe to
edit, save, build, move, or delete: later NeoSetzer updates do not overwrite
it. Build \texttt{main.tex} first, then keep this document open while you
explore the source and the PDF preview.

\subsection{Choose a starting point}
\label{sec:getting-started:starting-point}

The welcome screen offers several paths into a document. Create a blank LaTeX
file when you want a minimal starting point, or use the \textbf{Document
Wizard} to choose a document class, packages, title information, and other
sensible defaults. The wizard also keeps its small configuration presets
separate from \textbf{user document templates}: a user document template is a
snapshot of a complete LaTeX source document that can be saved, previewed, and
used as a starting point later.

This example project is deliberately another starting point. It is useful when
you prefer to learn by changing a working document rather than assembling every
piece from scratch.

\subsection{Build a first PDF}
\label{sec:getting-started:first-build}

The Magic Comment at the beginning of \texttt{main.tex} selects
\texttt{pdflatex}. NeoSetzer reads this small source-level instruction and
uses it only when the configured program is supported. The document keeps
PDFLaTeX as a portable baseline; you can change the configured build system in
\menu{Preferences} when your project needs XeLaTeX, LuaLaTeX, Tectonic, or
another available builder.

Use \textbf{Save and Build} after changing one sentence. The default shortcut
is commonly F5, but keyboard shortcuts are configurable, so the authoritative
list is \menu{Preferences \textrightarrow{} Keyboard Shortcuts}. The preview
shows the generated PDF, and the build log identifies errors, warnings, and bad
boxes with links back to source locations.

\subsection{A small task to try}
\label{sec:getting-started:task}

Change the title in \texttt{main.tex}, save, and build again. Then open the
Document Structure sidebar and select this section. It should take you to the
corresponding source heading. This short loop---edit, build, preview, and
navigate---is the foundation for the rest of the project.
